% 方法论章节

\section{方法论}

\subsection{符号定义}

为保持全文符号一致性，本文采用以下统一符号体系：

\begin{table}[H]
\centering
\caption{全文统一符号定义。包含数据集、特征向量、类别标签、模型预测函数、分类阈值、混淆矩阵元素和业务成本等核心符号。}
\label{tab:notation}
\begin{tabular}{cl}
\toprule
\textbf{符号} & \textbf{含义} \\
\midrule
$\mathcal{D}$ & 完整数据集 $\mathcal{D} = \{(\mathbf{x}_i, y_i)\}_{i=1}^{N}$ \\
$\mathbf{x}_i \in \mathbb{R}^d$ & 第$i$个样本的特征向量（$d=29$） \\
$y_i \in \{0, 1\}$ & 类别标签（$0$=正常交易，$1$=欺诈交易） \\
$\mathcal{D}_{\text{train}}, \mathcal{D}_{\text{val}}, \mathcal{D}_{\text{test}}$ & 训练集、验证集、测试集 \\
$N$ & 样本总数 \\
$d$ & 特征维度 \\
$f(\cdot)$ & 模型预测函数，$f: \mathbb{R}^d \to [0, 1]$ \\
$\hat{y}_i = f(\mathbf{x}_i)$ & 模型对样本$i$的预测概率 \\
$\tau$ & 分类阈值，$\hat{y}_i \geq \tau$ 判定为欺诈 \\
TP, FP, FN, TN & 真正例、假正例、假负例、真负例 \\
$\text{Cost}$ & 业务成本，$\text{Cost} = C_{FP} \cdot \text{FP} + C_{FN} \cdot \text{FN}$ \\
$C_{FP}, C_{FN}$ & 误报成本、漏报成本（$C_{FP}=1, C_{FN}=10$） \\
$\phi_j$ & 特征$j$的Shapley值 \\
\bottomrule
\end{tabular}
\end{table}

\subsection{数据来源与探索性分析}

本文采用Kaggle平台上的Credit Card Fraud Detection Dataset 2023数据集\cite{kaggle2023fraud}。该数据集包含信用卡交易记录，原始数据超过56.8万条，且正负类接近1:1平衡。然而，真实信用卡欺诈的比率通常仅为0.1\%--1\%，为模拟这一现实场景，本文对多数类（正常交易）进行下采样并限制最大样本数为150,000，同时从少数类（欺诈交易）中随机抽取750条，使最终欺诈比例约为0.5\%（正负比约1:200）。数据已通过PCA进行部分匿名化处理，但其规模（处理后150,750条记录）为深度学习模型提供了充足的训练样本。

\subsubsection{数据集特征}

数据集共包含31个特征列：28个匿名化PCA特征（V1-V28）、交易金额（Amount）和类别标签（Class，0=正常，1=欺诈）。初步探索性分析显示：

\begin{itemize}
    \item 总交易数：150,750（150,000正常 + 750欺诈）
    \item 欺诈交易占比：0.5\%（严重不平衡，正负比约1:200）
    \item 无缺失值，所有特征均为数值型
    \item Amount特征呈现典型的长尾分布
    \item 部分V特征在欺诈和正常交易间呈现显著分布差异
\end{itemize}

图\ref{fig:class_balance}呈现了经不平衡处理后的完整数据集类别分布：正常交易样本数为150,000，欺诈交易样本数仅为750，正负类比例约为200:1。这种极端的类别不平衡是信用卡欺诈检测任务的核心挑战——模型在训练过程中极易被多数类主导，导致对少数类的识别能力严重退化。仅0.5\%的欺诈比例意味着，即使一个"将所有交易预测为正常"的朴素分类器也能达到99.5\%的准确率——这正是准确率（Accuracy）在不平衡场景下完全失效的原因，也解释了本文为何采用AUPRC而非准确率作为主要评估指标。

图\ref{fig:amount_dist}以核密度估计对比了两类交易的金额分布。蓝色曲线（正常交易，$n=150,000$）和橙色曲线（欺诈交易，$n=750$）均呈现明显的右偏（长尾）分布，但存在两个关键差异：第一，欺诈交易的KDE峰值位置略高于正常交易，表明欺诈者倾向于选择中等偏高的交易金额以在"隐匿性"与"收益"之间取得平衡；第二，在高金额区间（$>15,000$标准化单位），欺诈交易的密度相对更高，这与银行业务经验一致——高价值交易是欺诈者的优先目标。然而，两类分布在整体上存在大量重叠区域，意味着仅凭交易金额远不足以区分欺诈与正常交易，必须依赖多维特征联合判别。

\begin{figure}[htbp]
    \centering
    \includegraphics[width=0.75\textwidth]{class_balance.pdf}
    \caption{经不平衡处理后的数据集类别分布。正常交易150,000条，欺诈交易750条，正负比约200:1，准确率在此场景下完全失效。}
    \label{fig:class_balance}
\end{figure}

\begin{figure}[htbp]
    \centering
    \includegraphics[width=0.78\textwidth]{amount_distribution.pdf}
    \caption{正常交易与欺诈交易的金额核密度估计曲线。蓝色为正常交易，橙色为欺诈交易，两类均呈右偏分布且存在大量重叠，欺诈交易在高金额区间密度略高。}
    \label{fig:amount_dist}
\end{figure}

上述EDA结果表明，任何有效的欺诈检测模型必须：（1）采用针对性的不平衡处理策略（如Focal Loss、SMOTE或类别加权），因为0.5\%的欺诈率使标准分类器失效；（2）依赖多特征联合判别而非单一金额阈值，因为两类金额分布存在大量重叠；（3）在评估时优先使用对类别不平衡敏感的指标（如AUPRC而非ROC-AUC或准确率）。

\subsection{数据预处理}

数据预处理流水线包含以下步骤：

\begin{enumerate}
    \item \textbf{特征选择}：移除ID列（不携带预测信息）。
    \item \textbf{异常值处理}：对Amount特征采用IQR方法进行异常值截断（仅在训练集上计算上下界，避免数据泄露）。V1-V28为PCA变换后的特征，其分布为原始特征的线性组合，不具有传统意义上的"异常值"含义，因此不做截断处理。
    \item \textbf{标准化}：使用StandardScaler对所有数值特征进行标准化（拟合于训练集，应用于验证集和测试集）。
    \item \textbf{数据集划分}：按70\%/15\%/15\%进行分层划分（Stratified Split），确保各个子集中欺诈比例一致。
\end{enumerate}

\subsection{模型架构}

本文实现并比较了两种异构深度神经网络架构与XGBoost传统基线。在此基础上，引入Tree-Supervised知识蒸馏方案（第三节），将XGBoost的树归纳偏置迁移至FT-Transformer：

\subsubsection{XGBoost基线}

XGBoost（eXtreme Gradient Boosting）\cite{chen2016xgboost}作为传统表格数据的强基线，采用梯度提升树框架。本文使用300轮提升、最大深度8、学习率0.05，并通过scale\_pos\_weight参数根据训练集正负比自动调节类别权重。为充分利用数据，XGBoost在训练集和验证集的合并数据上训练。

% =============================================================================
% FT-Transformer 原理深度讲解
% =============================================================================

\subsection{FT-Transformer：特征令牌化与自注意力机制}

FT-Transformer（Feature Tokenizer + Transformer）由Gorishniy等人\cite{gorishniy2021revisiting}于2021年提出，其核心思想是将表格数据的每个数值特征视为一个独立的"词元"（Token），通过Transformer编码器的自注意力机制学习特征间的全局交互。以下从四个层面逐层拆解其工作原理。

\subsubsection{整体架构}

FT-Transformer将表格分类问题转化为类似NLP的序列分类问题：每条数据记录被表示为一个特征令牌序列，模型通过在所有特征上计算成对注意力来捕获特征间关系。图\ref{fig:ftt_arch}以流程图呈现了从原始特征到最终预测的完整数据流。

\begin{figure}[htbp]
\centering
\begin{tikzpicture}[
    node distance=0.5cm,
    box/.style={rectangle, draw=black!70, rounded corners=3pt, minimum width=5.0cm, minimum height=0.62cm, align=center, font=\footnotesize},
    smallbox/.style={rectangle, draw=black!50, rounded corners=2pt, minimum width=2.2cm, minimum height=0.55cm, align=center, font=\scriptsize},
    arrow/.style={-{Stealth[scale=1.0]}, thick, black!60},
]
    % Main flow
    \node[box, fill=blue!5] (input) {输入 $\mathbf{x} = [x_1, \dots, x_{29}] \in \mathbb{R}^{29}$};
    \node[box, fill=blue!8, below=of input] (tokenizer) {Feature Tokenizer: $T_j = x_j \cdot \bm{W}_j + \bm{b}_j$ \quad ($\bm{W}_j,\bm{b}_j \in \mathbb{R}^{64}$)};
    \node[box, fill=green!5, below=of tokenizer] (cls) {前置 [CLS]: $\mathbf{Z} = [T_{\mathrm{CLS}}, T_1, \dots, T_{29}] \in \mathbb{R}^{30 \times 64}$};
    \node[box, fill=orange!8, below=of cls] (encoder) {Transformer Encoder $\times 3$\\(多头自注意力 + FFN + 残差连接 + LayerNorm)};
    \node[box, fill=green!5, below=of encoder] (extract) {提取 [CLS] 表示: $\mathbf{h}_{\mathrm{CLS}}^{(3)} \in \mathbb{R}^{64}$};
    \node[box, fill=red!5, below=of extract] (head) {LayerNorm $\to$ Linear($64{\to}1$) $\to$ Sigmoid};
    \node[box, fill=red!8, below=of head] (output) {输出: $\hat{y} \in [0, 1]$ （欺诈概率）};

    % Arrows
    \draw[arrow] (input) -- (tokenizer);
    \draw[arrow] (tokenizer) -- (cls);
    \draw[arrow] (cls) -- (encoder);
    \draw[arrow] (encoder) -- (extract);
    \draw[arrow] (extract) -- (head);
    \draw[arrow] (head) -- (output);

    % Encoder detail callout
    \node[smallbox, fill=orange!15, right=2.2cm of encoder] (attn) {Multi-Head\\Self-Attention};
    \node[smallbox, fill=orange!15, right=2.2cm of attn] (ffn) {Feed-Forward\\Network};
    \draw[arrow, dashed, orange!70] (encoder.east) -- ++(0.4,0) |- (attn.west);
    \draw[arrow, dashed, orange!70] (encoder.east) -- ++(0.4,0) |- (ffn.west);
\end{tikzpicture}
\caption{FT-Transformer架构流程图。29个数值特征经Feature Tokenizer映射为64维嵌入向量，前置可学习的[CLS] Token后送入3层Transformer编码器。[CLS]的最终输出经LayerNorm和线性层得到欺诈概率。右侧虚线框展示编码器内部两大核心组件。}
\label{fig:ftt_arch}
\end{figure}

\subsubsection{第一层：Feature Tokenizer——为什么数值需要"嵌入"}

传统做法将表格数据的一行直接输入MLP（$\mathbb{R}^{29} \to \text{Linear}(29, 64)$），这等价于对所有特征使用同一个投影矩阵。FT-Transformer的关键创新在于为每个特征$j$分配独立的可学习参数$\bm{W}_j$和$\bm{b}_j$：

\begin{equation}
T_j = x_j \cdot \bm{W}_j + \bm{b}_j \in \mathbb{R}^{d_{model}}
\label{eq:tokenize}
\end{equation}

其中$\bm{W}_j$的作用不是线性变换（因为$x_j$是标量），而是\textbf{特征专属的嵌入向量}——将标量特征值沿特定方向"拉伸"为$d_{model}$维向量。理解这一设计的关键在于：

\begin{itemize}
    \item \textbf{异构特征需要异构空间}：V14（强判别性PCA成分）和V25（近噪声成分）的变化量$\Delta x$对最终决策的影响截然不同。若所有特征共享同一投影矩阵，模型将无法区分"重要的0.1"和"不重要的0.1"。Feature Tokenizer使每个特征拥有独立的"语义通道"，其嵌入向量的方向编码了该特征的身份（identity），模长编码了其数值强度。
    \item \textbf{为注意力提供可区分的Key/Query}：自注意力通过Query-Key内积$\langle \bm{Q}_i, \bm{K}_j \rangle$度量特征$i$对特征$j$的"关注程度"。若特征嵌入是同一线性变换的产物，则所有特征将位于同一低维子空间，Query-Key内积丧失区分度——这正是本文在PCA匿名化数据上观察到的"注意力均匀塌缩"的根本原因。
\end{itemize}

在本文的实验中，$d_{model}=64$意味着每个标量特征被映射为一个64维向量，29个特征共计产生$29 \times (64 + 64) = 3,\!712$个专属参数。

\subsubsection{第二层：[CLS] Token——从序列到分类的桥梁}

借鉴BERT \cite{devlin2018bert}的设计，FT-Transformer在所有特征令牌前添加一个可学习的[CLS] Token $T_{\mathrm{CLS}} \in \mathbb{R}^{d_{model}}$，形成长度为30的序列：

\begin{equation}
\mathbf{Z}_0 = [T_{\mathrm{CLS}}, T_1, T_2, \dots, T_{29}] \in \mathbb{R}^{30 \times 64}
\end{equation}

[CLS] Token的设计哲学是\textbf{信息汇聚}：它自身不携带任何特征信息（随机初始化），但在经过多层自注意力后，它通过Query机制从所有特征令牌中"抽取"了对分类最有用的信息。可以理解为——[CLS]是模型学会的"最佳问题"（optimal query），而每个特征令牌提供的是该特征的"回答"（key-value pair）。最终，仅[CLS]的输出被送入分类头，其余29个特征令牌的输出被丢弃。

\subsubsection{第三层：Multi-Head Self-Attention——特征间的全局对话}

这是FT-Transformer的核心计算引擎。给定输入序列$\mathbf{Z} \in \mathbb{R}^{S \times d}$（$S=30$），单头自注意力执行以下操作：

\textbf{Step 1：线性投影。}将每个令牌分别投影为Query、Key、Value三个向量：
\begin{equation}
\mathbf{Q} = \mathbf{Z}\mathbf{W}_Q,\quad \mathbf{K} = \mathbf{Z}\mathbf{W}_K,\quad \mathbf{V} = \mathbf{Z}\mathbf{W}_V
\end{equation}
其中$\mathbf{W}_Q, \mathbf{W}_K, \mathbf{W}_V \in \mathbb{R}^{d \times d_k}$，$d_k = d_{model} / n_{heads} = 8$。

\textbf{Step 2：Scaled Dot-Product Attention。}计算每对令牌间的注意力分数：
\begin{equation}
\mathrm{Attention}(\mathbf{Q}, \mathbf{K}, \mathbf{V}) = \mathrm{Sparsemax}\!\left(\frac{\mathbf{Q}\mathbf{K}^\top}{\sqrt{d_k}}\right) \mathbf{V}
\label{eq:attention}
\end{equation}

除以$\sqrt{d_k}$的目的是防止内积过大导致Sparsemax/Softmax饱和（梯度消失）。30个令牌的注意力矩阵为$\mathbb{R}^{30 \times 30}$——每个特征对其他29个特征（包括[CLS]和自身）分配一个注意力权重。

\textbf{Step 3：Multi-Head 整合。}并行执行$n_{heads}=8$个独立的注意力计算（每头使用不同的$\mathbf{W}_Q, \mathbf{W}_K, \mathbf{W}_V$），将各头输出拼接后线性投影：

\begin{equation}
\mathrm{MultiHead}(\mathbf{Z}) = \mathrm{Concat}(\mathrm{head}_1, \dots, \mathrm{head}_8) \mathbf{W}_O
\end{equation}

多头设计使模型能在不同的表示子空间中同时关注不同类型的关系——例如，一个头可能聚焦于V14-V12的相关性，另一个头关注V4-V17的交互。然而，本文的实验发现：在PCA匿名化特征上，各头输出的注意力分布高度相似，表明PCA破坏了供多头分化的语义多样性。

\textbf{Step 4：Transformer Encoder Block。}每个编码器层包含两个子层，每子层均有残差连接和LayerNorm：

\begin{equation}
\begin{aligned}
\mathbf{Z}' &= \mathrm{LayerNorm}\big(\mathbf{Z} + \mathrm{MultiHead}(\mathbf{Z})\big) \\
\mathbf{Z}'' &= \mathrm{LayerNorm}\big(\mathbf{Z}' + \mathrm{FFN}(\mathbf{Z}')\big)
\end{aligned}
\end{equation}

其中FFN为两层MLP：$\mathrm{FFN}(\mathbf{x}) = \mathrm{GELU}(\mathbf{x}\mathbf{W}_1 + \mathbf{b}_1)\mathbf{W}_2 + \mathbf{b}_2$，中间维度$d_{ffn}=256$提供4倍于$d_{model}$的容量以进行非线性变换。残差连接确保即使注意力学不到有用信息（如均匀塌缩），梯度仍能通过恒等路径回传，避免深层网络的退化。

\subsubsection{第四层：Sparsemax——本文的关键改进}

原始FT-Transformer使用Softmax归一化注意力权重：
\begin{equation}
\mathrm{Softmax}(\mathbf{z})_i = \frac{e^{z_i}}{\sum_j e^{z_j}}
\end{equation}

Softmax存在两个固有问题：（1）它永远输出全非零的概率分布（$e^{z_i} > 0$恒成立），无法产生真正的稀疏解；（2）当所有$z_i$接近时（如在PCA空间中），输出趋于均匀分布$1/d$。

本文采用Sparsemax \cite{martins2016sparsemax}替代Softmax。Sparsemax将logits向量$\mathbf{z}$投影到概率单纯形（probability simplex）上，产生\textbf{可能包含零值}的稀疏概率分布：

\begin{equation}
\mathrm{Sparsemax}(\mathbf{z}) = \arg\min_{\mathbf{p} \in \Delta^{d-1}} \|\mathbf{p} - \mathbf{z}\|_2^2
\end{equation}

其闭式解为$\mathrm{Sparsemax}(\mathbf{z})_i = \max(z_i - \tau(\mathbf{z}), 0)$，其中$\tau(\mathbf{z})$是确保$\sum_i \max(z_i - \tau, 0) = 1$的唯一阈值。算法复杂度为$O(d \log d)$，不增加计算负担。

\textbf{核心差异}：Softmax对低logits值赋予微量但非零的概率（如0.001），Sparsemax直接将其设为零。在表格数据语境中，这意味着特征$j$完全可以被排除在特征$i$的"注意力视野"之外——模型被强制做出选择。本实验证实，Sparsemax将FT-Transformer的AUPRC从0.8738提升至0.9179（+4.4pp），ECE从0.2328降至0.0120（改善约20倍）。

需要指出，本文未分离Sparsemax的稀疏化效应与其隐式的温度/熵调节效应——Sparsemax的投影操作（$\ell_2$范数最小化至单纯形）在降低输出熵的同时引入稀疏性，两者可能独立贡献于AUPRC和ECE的改善。Softmax+温度缩放（temperature scaling）可能以不同机制实现类似的熵调节效果，这一替代路径留待未来验证。

\subsubsection{本文配置与训练细节}

\begin{itemize}
    \item $d_{model}=64, n_{heads}=8, n_{layers}=3, d_{ffn}=256$
    \item 总参数量：153,921（合理范围，避免过拟合15万训练样本）
    \item 优化器：AdamW（lr=$5 \times 10^{-4}$, weight\_decay=$10^{-4}$）
    \item 损失函数：Focal Loss（$\alpha=0.25, \gamma=2.0$）——$\alpha=0.25$被选为Sparsemax的搭配，因稀疏注意力需要均衡的正负类梯度信号
    \item 早停：验证集AUPRC连续15轮不提升
\end{itemize}

% =============================================================================
% DeepFM 原理深度讲解
% =============================================================================

\subsection{DeepFM：显式特征交互与因子分解机}

DeepFM由Guo等人\cite{guo2017deepfm}于2017年提出，最初应用于点击率（CTR）预测。其核心思想是将因子分解机（Factorization Machine, FM）的显式二阶特征交互与深度神经网络的高阶表示学习融合于一个端到端框架中。与FT-Transformer的"隐式"交互（通过注意力权重间接建模）不同，DeepFM采用"显式"交互——直接计算每对特征的向量内积。

\subsubsection{整体架构}

DeepFM由三个组件构成，共享底层的特征嵌入层。图\ref{fig:deepfm_arch}展示了从前向传播到最终预测的完整数据流。

\begin{figure}[htbp]
\centering
\begin{tikzpicture}[
    node distance=0.5cm,
    box/.style={rectangle, draw=black!70, rounded corners=3pt, minimum width=5.4cm, minimum height=0.6cm, align=center, font=\footnotesize},
    branch/.style={rectangle, draw=black!50, rounded corners=2pt, minimum width=2.6cm, minimum height=0.58cm, align=center, font=\scriptsize},
    arrow/.style={-{Stealth[scale=1.0]}, thick, black!60},
    merge/.style={circle, draw=black!70, fill=black!5, minimum size=0.55cm, font=\small},
]
    % Main flow
    \node[box, fill=blue!5] (input) {输入 $\mathbf{x} = [x_1, \dots, x_{29}] \in \mathbb{R}^{29}$};
    \node[box, fill=green!8, below=of input] (embed) {共享嵌入: $\mathbf{e}_i = x_i \cdot \mathbf{v}_i \in \mathbb{R}^{16}$ ($i = 1,\dots,29$)};
    \node[box, fill=orange!5, below=0.15cm of embed] (split) {分两路：FM路径 + Deep路径（共享嵌入参数）};

    % FM branch
    \node[branch, fill=yellow!8, below left=0.6cm and -1.6cm of split] (linear) {一阶线性: $\sum_{i} w_i x_i$};
    \node[branch, fill=yellow!8, below=0.3cm of linear] (second) {二阶交互: $\sum_{i<j}\langle\mathbf{v}_i,\mathbf{v}_j\rangle x_i x_j$};

    % Deep branch
    \node[branch, fill=cyan!8, below right=0.6cm and -1.6cm of split] (dense1) {Dense(464$\to$256) + ReLU};
    \node[branch, fill=cyan!8, below=0.3cm of dense1] (dense2) {Dense(256$\to$128) + ReLU};
    \node[branch, fill=cyan!8, below=0.3cm of dense2] (dense3) {Dense(128$\to$64) + ReLU};

    % Merge
    \node[merge, below=2.0cm of split] (merge) {$\Sigma$};
    \node[box, fill=red!5, below=0.35cm of merge] (sigmoid) {$\sigma\!\left(w_0 + \text{FM} + \text{Deep}\right)$};
    \node[box, fill=red!8, below=0.35cm of sigmoid] (output) {输出: $\hat{y} \in [0, 1]$};

    % Arrows from input
    \draw[arrow] (input) -- (embed);
    \draw[arrow] (embed) -- (split);

    % FM branch arrows
    \draw[arrow] (split.south) -- ++(0, -0.25) -| (linear.north);
    \draw[arrow] (linear) -- (second);
    \draw[arrow] (second.south) -- ++(0, -0.2) -| (merge.west);

    % Deep branch arrows
    \draw[arrow] (split.south) -- ++(0, -0.25) -| (dense1.north);
    \draw[arrow] (dense1) -- (dense2);
    \draw[arrow] (dense2) -- (dense3);
    \draw[arrow] (dense3.south) -- ++(0, -0.2) -| (merge.east);

    % Merge to output
    \draw[arrow] (merge) -- (sigmoid);
    \draw[arrow] (sigmoid) -- (output);

    % Annotation
    \node[font=\scriptsize\itshape, black!45, right=1.0cm of second] {FM: $O(DK)$};
    \node[font=\scriptsize\itshape, black!45, right=1.0cm of dense3] {Deep: MLP};
\end{tikzpicture}
\caption{DeepFM架构流程图。29个数值特征经共享嵌入层映射为16维向量，分两路并行计算：FM路径执行显式一阶和二阶特征交互，Deep路径通过三层MLP学习高阶表示。两路输出与全局偏置$w_0$求和后经Sigmoid输出。共享嵌入是DeepFM的核心设计——FM的交互质量直接影响Deep路径的初始表示。}
\label{fig:deepfm_arch}
\end{figure}

\subsubsection{第一层：共享嵌入}

每个特征$x_i$（标量）被映射为一个$K$维嵌入向量（本文$K=16$）：

\begin{equation}
\mathbf{e}_i = x_i \cdot \mathbf{v}_i \in \mathbb{R}^{K}
\label{eq:deepfm_embed}
\end{equation}

其中$\mathbf{v}_i \in \mathbb{R}^{K}$是特征$i$的专属嵌入向量。所有这些$\mathbf{v}_i$在FM组件和Deep组件之间\textbf{共享}——这意味着FM学习到的成对关系会通过反向传播同时优化Deep路径使用的嵌入表示。共享机制是DeepFM超越"FM+MLP串联"的关键：两个组件的训练信号在嵌入层交汇，FM提供的结构化交互信号充当Deep路径的正则化先验。

\subsubsection{第二层：FM组件}

FM组件由两部分组成：

\textbf{一阶部分（线性）：}
\begin{equation}
y_{\mathrm{linear}} = w_0 + \sum_{i=1}^{D} w_i x_i
\end{equation}

这部分等价于逻辑回归——每个特征有独立的线性权重$w_i$，直接对预测值做加减。$w_0$是全局偏置。

\textbf{二阶部分（成对交互）：}
\begin{equation}
y_{\mathrm{FM}} = \sum_{i=1}^{D} \sum_{j=i+1}^{D} \langle \mathbf{v}_i, \mathbf{v}_j \rangle \cdot x_i x_j
\label{eq:fm_interaction}
\end{equation}

这是FM的核心。对每一对特征$(i, j)$，其交互强度由两个因素共同决定：
\begin{enumerate}
    \item \textbf{数值乘积} $x_i \cdot x_j$：如果两个特征同时偏高（如V14和V12均异常高），乘积大→交互贡献大。
    \item \textbf{嵌入相似度} $\langle \mathbf{v}_i, \mathbf{v}_j \rangle$：两个特征的嵌入向量越"对齐"（内积越大），它们对预测的协同效应越强。
\end{enumerate}

\textbf{关键数学技巧——$O(D^2K) \to O(DK)$：}直接计算公式\eqref{eq:fm_interaction}需要对$D(D-1)/2$对特征各计算一次$K$维内积，复杂度$O(D^2K)$。通过以下代数变换可降至$O(DK)$：

\begin{equation}
\begin{aligned}
\sum_{i=1}^{D}\sum_{j=i+1}^{D} \langle \mathbf{v}_i, \mathbf{v}_j \rangle x_i x_j
&= \frac{1}{2}\left(\sum_{i=1}^{D}\sum_{j=1}^{D} \langle \mathbf{v}_i, \mathbf{v}_j \rangle x_i x_j - \sum_{i=1}^{D} \langle \mathbf{v}_i, \mathbf{v}_i \rangle x_i^2 \right) \\
&= \frac{1}{2}\left( \sum_{f=1}^{K}\left(\sum_{i=1}^{D} v_{i,f} x_i\right)^2 - \sum_{i=1}^{D}\sum_{f=1}^{K} v_{i,f}^2 x_i^2 \right)
\end{aligned}
\label{eq:fm_trick}
\end{equation}

式\eqref{eq:fm_trick}的关键在于将双重求和$\sum_i\sum_j$转化为先对$i$求和再平方的结构——这只需要$O(DK)$次乘法和加法，而非$O(D^2K)$次。当$D=29$时，$O(DK)=464$次操作 vs. $O(D^2K)=13,456$次操作，加速约29倍。这一变换使得FM的大规模特征交互在计算上可行。

\subsubsection{第三层：Deep组件}

Deep组件是一个标准的前馈神经网络，接收所有特征嵌入向量的拼接（flatten）作为输入：

\begin{equation}
\mathbf{h}_0 = \mathrm{Concat}(\mathbf{e}_1, \mathbf{e}_2, \dots, \mathbf{e}_{29}) \in \mathbb{R}^{D \times K = 464}
\end{equation}

随后通过三层全连接网络：

\begin{equation}
\begin{aligned}
\mathbf{h}_1 &= \mathrm{ReLU}(\mathbf{W}_1 \mathbf{h}_0 + \mathbf{b}_1), \quad \mathbf{W}_1 \in \mathbb{R}^{256 \times 464} \\
\mathbf{h}_2 &= \mathrm{ReLU}(\mathbf{W}_2 \mathbf{h}_1 + \mathbf{b}_2), \quad \mathbf{W}_2 \in \mathbb{R}^{128 \times 256} \\
\mathbf{h}_3 &= \mathrm{ReLU}(\mathbf{W}_3 \mathbf{h}_2 + \mathbf{b}_3), \quad \mathbf{W}_3 \in \mathbb{R}^{64 \times 128}
\end{aligned}
\end{equation}

每层后应用Dropout（$p=0.3$）防止过拟合。Deep组件可以学习任意阶的特征交互——例如，V14$\times$V12$\times$V4的三阶交互可以通过两层ReLU的非线性组合隐式建模。

\subsubsection{第四层：最终融合}

三路输出求和后经Sigmoid得到最终预测：

\begin{equation}
\hat{y} = \sigma\Big(w_0 + \underbrace{\sum_{i=1}^{29} w_i x_i}_{\text{一阶线性}} + \underbrace{\frac{1}{2}\sum_{f=1}^{16}\Big[\big(\sum_i v_{i,f}x_i\big)^2 - \sum_i v_{i,f}^2 x_i^2\Big]}_{\text{FM二阶交互}} + \underbrace{\mathbf{W}_{\mathrm{out}} \cdot \mathbf{h}_3}_{\text{Deep高阶}} \Big)
\end{equation}

为什么使用"求和"而非"拼接"融合？FM和Deep组件共享嵌入层，它们的输出已经位于相关的语义空间。求和是一种参数无关的融合方式，允许模型通过反向传播自动分配两个组件的相对贡献——如果FM的二阶交互足够，Deep输出趋近于零；如果高阶模式更关键，Deep输出主导预测。

\subsubsection{DeepFM在PCA空间中的优势}

本文的关键实验发现——DeepFM（AUPRC=0.9336）优于FT-Transformer（0.9179）——可从架构设计中得到解释：

\begin{enumerate}
    \item \textbf{不依赖特征"身份"的全局比较}：FT-Transformer的注意力通过Query-Key内积评估"特征$i$对特征$j$有多重要"，这要求特征嵌入空间具有可区分的几何结构。PCA匿名化抹去了特征的原始语义，使各嵌入向量在空间中高度同质化，内积趋同→注意力均匀。DeepFM的$\langle \mathbf{v}_i, \mathbf{v}_j \rangle$内积只是两个嵌入向量的"对齐程度"的静态度量，不依赖Query-Key的动态上下文——它在训练前就确定了"V14和V12的交互值得关注"。
    \item \textbf{交互的直接性}：FM的二阶交互是$x_i \cdot x_j \cdot \langle \mathbf{v}_i, \mathbf{v}_j \rangle$的简单三因子乘积，梯度路径短且稳定。FT-Transformer的交互需要经过Softmax/Sparsemax归一化+Value加权求和+残差连接，梯度需要穿越多个非线性层。
    \item \textbf{知识蒸馏的兼容性}：XGBoost的分裂决策本质上是轴对齐的——每次分裂沿单一特征维度判断$x_i > t$。DeepFM的一阶线性项$w_i x_i$直接对应这种轴对齐模式，软标签蒸馏的信号得以无损注入。而FT-Transformer的注意力无视特征轴，蒸馏信号在注意力矩阵中经过非线性变换后被稀释。
\end{enumerate}

\subsubsection{本文配置与训练细节}

\begin{itemize}
    \item $K=16$（嵌入维度），Deep layers=[256, 128, 64]，Dropout=0.3
    \item 总参数量：161,648（略高于FT-Transformer的153,921）
    \item 优化器：AdamW（lr=$10^{-3}$, weight\_decay=$10^{-4}$）
    \item 损失函数：Focal Loss（$\alpha=0.995$即类别反比, $\gamma=2.0$）——DeepFM使用强类别加权以弥补其缺乏内置稀疏注意力（Sparsemax）的劣势
    \item 蒸馏配置（DeepFM+KD）：$T=4.0, \alpha_{\mathrm{KD}}=0.7$——软标签权重占30\%，Focal Loss占70\%
\end{itemize}


\subsection{训练策略}

\subsubsection{Focal Loss}

本文采用Focal Loss作为主要损失函数\cite{lin2017focal}：

\begin{equation}
\text{FL}(p_t) = -\alpha_t (1 - p_t)^\gamma \log(p_t)
\end{equation}

其中 $p_t = p$ 当 $y=1$ 否则 $p_t = 1-p$，$\alpha_t = \alpha$ 当 $y=1$ 否则 $\alpha_t = 1-\alpha$。

两个关键超参数：
\begin{itemize}
    \item $\gamma$（聚焦参数）：控制对易分类样本的"惩罚"力度。$\gamma=0$退化为标准交叉熵，$\gamma$越大，模型越关注难分类样本。
    \item $\alpha$（平衡参数）：调节正负类之间的权重。对于FT-Transformer，因其Sparsemax注意力需要更均衡的梯度信号，使用默认值$\alpha=0.25$；对于DeepFM，设$\alpha = 1 - \text{正类比例}$，使少数类获得更高的损失权重。
\end{itemize}

本文将Focal Loss作为主要损失函数用于所有深度模型训练。

\subsubsection{优化设置}

所有神经网络模型使用AdamW优化器（$\beta_1=0.9, \beta_2=0.999$），配合ReduceLROnPlateau学习率调度（当验证集AUPRC连续8轮无提升时，学习率减半）。训练采用以下策略：
\begin{itemize}
    \item 批量大小：512
    \item 最大训练轮数：100
    \item 早停策略：验证集AUPRC连续15轮无提升时停止
    \item 梯度裁剪：max\_norm=1.0
    \item 最佳模型保存：基于验证集AUPRC
\end{itemize}

\subsection{评估指标}

\subsubsection{主要指标：AUPRC}

鉴于类别高度不平衡，AUPRC（Area Under Precision-Recall Curve）是比ROC-AUC更合适的评估指标\cite{saito2015precision}。ROC-AUC在不平衡数据上可能给出过于乐观的评估，因为真负类（True Negative）占主导地位。PR曲线聚焦于正类（欺诈）的预测质量，对样本分布变化更为敏感。

\subsubsection{次要指标}

\begin{itemize}
    \item ROC-AUC：便于与文献中的现有工作对比
    \item F1-Score：在最佳阈值（基于验证集F1最大化）处计算
    \item Precision \& Recall：分别衡量误报率和漏报率
\end{itemize}

\subsubsection{商业成本指标}

从银行业务角度出发，误报（将正常交易标记为欺诈）和漏报（未检出欺诈交易）的成本截然不同：

\begin{equation}
\text{Cost} = FP \times C_{FP} + FN \times C_{FN}, \quad C_{FP}=1, C_{FN}=10
\end{equation}

漏报（$FN$）的真实成本大约是误报（$FP$）的10倍——未发现的欺诈交易导致直接资金损失，而误报仅需额外的人工审核。
